\chapter{Requerimientos}\label{chapter03}

A continuación se detallan los requerimientos funcionales y no funcionales del sistema, junto con las historias de usuario que guían el desarrollo del producto.

\section{Historias de Usuario}

\begin{itemize}
	\item HU01 -- Como vendedor, quiero visualizar un ranking de productos en tendencia por categoría para decidir qué publicar.
	\item HU02 -- Como vendedor, quiero recibir una alerta cuando un producto que sigo empiece a crecer en tendencia, sin tener que entrar a revisar el dashboard.
	\item HU03 -- Como comprador, quiero ver el histórico de precio de un producto antes de comprarlo.
	\item HU04 -- Como usuario, quiero filtrar productos por categoría, rango de fecha y precio.
	\item HU05 -- Como administrador, quiero configurar la frecuencia de captura de datos desde la API de Mercado Libre (implementada mediante tareas programadas del sistema operativo: tres corridas diarias parametrizables).
	\item HU06 -- Como cliente, quiero crear una cuenta para acceder a las funciones del nivel pago del modelo freemium (historial completo, exportación y alertas).
	\item HU07 -- Como cliente, quiero seguir un producto puntual desde el ranking, para que el sistema me avise si empieza a crecer en vez de tener que revisarlo yo.
\end{itemize}

\section{Requerimientos Funcionales}

\begin{enumerate}
	\item RF01 -- El sistema debe capturar periódicamente datos de productos desde la API oficial de Mercado Libre.
	\item RF02 -- El sistema debe calcular un score de tendencia por producto en base a variables temporales.
	\item RF03 -- El sistema debe permitir filtrar y ordenar productos por categoría, score, variación de precio, rango de precio y rango de fecha.
	\item RF04 -- El sistema debe exponer los resultados mediante una API REST (Next.js API Routes).
	\item RF05 -- El dashboard debe mostrar gráficos de evolución temporal por producto y categoría.
	\item RF06 -- El sistema debe permitir exportar el listado de productos en tendencia (CSV).
	\item RF07 -- El sistema debe exponer, para cada producto, la probabilidad estimada por el modelo de Machine Learning supervisado (Sección~\ref{sec:ml-resultados}) de que su score siga creciendo.
	\item RF08 -- El sistema debe permitir a un usuario crear una cuenta y autenticarse, para diferenciar el nivel gratuito del nivel pago del modelo freemium (Sección~\ref{sec:modelo-negocio}).
	\item RF09 -- El sistema debe permitir a un cliente autenticado seguir un producto puntual y generar una alerta cuando su score cruce el umbral de crecimiento (HU07).
\end{enumerate}

\section{Requerimientos No Funcionales}

\begin{enumerate}
	\item RNF01 -- La captura de datos no debe exceder los límites de rate-limit de la API oficial de Mercado Libre.
	\item RNF02 -- El dashboard debe responder consultas en menos de 3 segundos sobre el dataset histórico. El modelo de datos incluye índices en PostgreSQL (productId, timestamp, categoryId) diseñados para sostener ese tiempo de respuesta a medida que el histórico escala.
	\item RNF03 -- El backend debe poder desplegarse en plataformas cloud con escalado gestionado. Actualmente se despliega en Render con build reproducible (generación de Prisma Client en postinstall).
	\item RNF04 -- Los datos históricos deben persistirse con respaldo (backup) periódico.
	\item RNF05 -- Las credenciales y tokens de acceso (OAuth2 de Mercado Libre) deben gestionarse exclusivamente del lado servidor, sin exponerse nunca al cliente, con renovación automática de tokens.
	\item RNF06 -- El panel administrativo (\texttt{/admin}) no debe ser accesible sin autenticación, dado que expone información operativa interna (logs de minería, estado de endpoints).
	\item RNF07 -- Las contraseñas de los usuarios deben almacenarse con hash seguro (nunca en texto plano ni reversible).
\end{enumerate}

\section{Matriz de trazabilidad}\label{sec:trazabilidad}

La siguiente matriz vincula cada requerimiento con su estado real de implementación al cierre de este hito y con la evidencia concreta que lo respalda, medida sobre el sistema en producción (ver evidencia cuantitativa detallada en la Sección~\ref{sec:evidencia-cuantitativa}).

\begin{longtable}{@{}>{\raggedright\arraybackslash}p{1.6cm} >{\raggedright\arraybackslash}p{2.6cm} >{\raggedright\arraybackslash}p{7.8cm} >{\raggedright\arraybackslash}p{2.1cm}@{}}
	\caption{Matriz de trazabilidad de requerimientos} \label{tab:trazabilidad} \\
	\toprule
	\textbf{Req.} & \textbf{Estado} & \textbf{Evidencia} & \textbf{Sección} \\
	\midrule
	\endfirsthead
	\toprule
	\textbf{Req.} & \textbf{Estado} & \textbf{Evidencia} & \textbf{Sección} \\
	\midrule
	\endhead
	RF01 & Implementado & 274 productos y 5717 apariciones capturados en 71 días de operación continua (14/7 al 23/9/2026), en $\sim$203 sesiones de minería estimadas. & Sec.~\ref{sec:evidencia-cuantitativa}, Fig.~\ref{fig:diagrama-componentes} \\
	RF02 & Implementado & 11573+ scores calculados y persistidos con desglose completo. Fórmula, pesos y ejemplo real documentados. Desde el 23/09/2026 el componente de estabilidad se calcula sobre el precio promedio entre vendedores activos (no el más barato); verificado con una corrida real de minería, 12 productos con datos de múltiples vendedores cargados. & Sec.~\ref{sec:score-formal} \\
	RF03 & Implementado & Filtro por categoría, precio (mín./máx.) y orden por score/variación en el ranking; filtro por rango de fecha (7/30 días/todo) en el gráfico de evolución. & Fig.~\ref{fig:demo-score} \\
	RF04 & Implementado & 5 endpoints REST en producción (\texttt{/api/trend-scores}, \texttt{/api/health}, \texttt{/api/status}, \texttt{/api/search}, \texttt{/api/trends}), más los de cuentas y alertas (\texttt{/api/account/*}, \texttt{/api/watchlist}, \texttt{/api/alerts}). & Sec.~\ref{sec:evidencia-cuantitativa} \\
	RF05 & Implementado & Gráfico de evolución del score de tendencia por producto y panel de tendencias desglosado por categoría, ambos en el dashboard. & Fig.~\ref{fig:demo-evolucion} \\
	RF06 & Implementado & Exportación a CSV del ranking filtrado, disponible en el dashboard (nivel pago, Sección~\ref{sec:modelo-negocio}). & Fig.~\ref{fig:demo-score} \\
	RF07 & Implementado & Columna ``Prob. ML'' en el ranking con la probabilidad calculada por el checkpoint del modelo supervisado, cargada para los 227 productos con score. & Sec.~\ref{sec:ml-resultados} \\
	RF08 & Implementado & Registro y login con contraseña hasheada (bcrypt) y sesión firmada (JWT); nivel gratuito vs. pago diferenciado en el dashboard. & Sec.~\ref{sec:modelo-negocio} \\
	RF09 & Implementado & Botón de seguir por producto (requiere cuenta) y bandeja de alertas en \texttt{/alerts}, generadas al cruzar el umbral de crecimiento. & -- \\
	RNF01 & Implementado & Sin bloqueos por rate-limit en 71 días de operación; los únicos 2 endpoints bloqueados lo están por política de Mercado Libre, no por exceso de solicitudes. & Sec.~\ref{sec:evidencia-cuantitativa} \\
	RNF02 & Implementado & Medido en producción el 21/09/2026, tras acotar la consulta a las últimas 2 mediciones por producto: \texttt{/api/trend-scores} responde en 1,7--2,9\,s en caliente, dentro del objetivo de 3\,s (antes: 6,3--7,6\,s). Detalle de la corrección y su verificación en la Sección~\ref{sec:pruebas}. & Sec.~\ref{sec:evidencia-cuantitativa}, \ref{sec:pruebas} \\
	RNF03 & Implementado & Desplegado en Render, con build reproducible (generación de Prisma Client en \texttt{postinstall}). & -- \\
	RNF04 & Implementado & Respaldo automático (\texttt{backup\_db.js}) ejecutado al final de cada corrida de minería. & -- \\
	RNF05 & Implementado & Tokens OAuth2 persistidos server-side (tabla \texttt{OAuthToken}), con renovación automática cuando faltan menos de 5 minutos para expirar. & -- \\
	RNF06 & Implementado & \texttt{/admin} protegido por un \emph{proxy} (\texttt{src/proxy.ts}) que exige sesión con rol administrador; antes de esta corrección era accesible sin ninguna autenticación. & -- \\
	RNF07 & Implementado & Contraseñas hasheadas con bcrypt (factor de costo 10) antes de persistirse; nunca se guarda ni se loguea la contraseña en texto plano. & -- \\
	\bottomrule
\end{longtable}

Las historias de usuario se consideran implementadas de forma directa cuando su(s) RF asociado(s) están implementados: HU01 (RF01--RF03), HU03 (RF05), HU04 (RF03) y HU05 (RF01). HU02 y HU07 quedan cubiertas por RF09 (seguir un producto y recibir una alerta in-app cuando cruza el umbral de crecimiento) -- con una limitación reconocida: la alerta hoy se ve en \texttt{/alerts} dentro de la aplicación, pero todavía no se envía por email o push, porque eso requiere dar de alta una cuenta en un proveedor externo de correo (Sección~\ref{sec:alcance}). HU06 queda cubierta por RF08.

\section{Metodología de desarrollo}\label{sec:metodologia}

El desarrollo del proyecto se llevó adelante bajo un enfoque ágil de tipo Kanban: ciclos cortos e incrementales de trabajo, integración continua hacia el entorno de producción y ciclos de feedback frecuentes con el tutor, en lugar de iteraciones de duración fija (\emph{sprints}) cerradas de antemano y definidas por adelantado. \textcite{BeckEtAl2001} definen los principios centrales de este tipo de metodologías -- entrega temprana y continua de software funcionando, e incorporación del cambio incluso en etapas avanzadas del desarrollo -- que se reflejan directamente en la forma de trabajo adoptada en este proyecto.

Esta elección se justifica por tres motivos concretos, verificables en el historial de control de versiones del proyecto:

\begin{itemize}
	\item \textbf{Incrementos pequeños y frecuentes}: 63 \emph{commits} entre el 21 de abril y el 31 de agosto de 2026 (aproximadamente cuatro meses y medio), cada uno acotado a un cambio puntual -- una corrección, una funcionalidad o un ajuste de documentación -- en lugar de entregas grandes y espaciadas.
	\item \textbf{Integración continua}: cada incremento se despliega automáticamente al entorno de producción en Render ante cada cambio en la rama principal (RNF03, Sección~\ref{sec:trazabilidad}), permitiendo validar el sistema en condiciones reales en vez de acumular trabajo sin probar.
	\item \textbf{Reorganización del trabajo en respuesta a feedback}: el plan de tareas se ajustó explícitamente en más de una ocasión a partir de devoluciones del tutor -- documentadas tras la entrega del 25\,\% y nuevamente tras la del 50\,\% --, incorporando esos cambios como parte del flujo normal de trabajo en lugar de tratarlos como una fase separada de corrección final.
\end{itemize}

Se descartó una metodología estructurada de tipo cascada porque el proyecto combina componentes con distinto grado de incertidumbre: el pipeline de captura y el dashboard tenían requerimientos claros desde el inicio, mientras que el diseño del score de tendencia y, más adelante, el del modelo de Machine Learning supervisado (Sección~\ref{sec:evolucion-tecnica}) se fueron ajustando a medida que se disponía de más datos reales -- algo que un enfoque secuencial con etapas cerradas de antemano no habría permitido sin retrabajo significativo.

\section{Viabilidad legal y técnica}

El proyecto no aborda una temática sensible en términos legales -- no procesa datos de salud, alimentación ni pone en riesgo la vida de los usuarios --: los datos de producto (precios, categorías, publicaciones) provienen exclusivamente de la API oficial de Mercado Libre (Sección~\ref{sec:tecnologias-utilizadas}) y son públicos, sin acceder a información privada de compradores o vendedores de Mercado Libre (Sección~\ref{sec:alcance}). Por ese motivo, no corresponde incluir un descargo de responsabilidad legal específico más allá del cumplimiento de los términos de uso de la API oficial, ya documentado como RNF05 (gestión segura de credenciales OAuth2 exclusivamente del lado servidor, Sección~\ref{sec:trazabilidad}).

Con la incorporación de cuentas de usuario (RF08, Sección~\ref{sec:modelo-negocio}) el sistema sí pasó a recolectar un dato personal propio, no de Mercado Libre: el email de quien se registra. El tratamiento se mantiene acotado a lo estrictamente necesario para el modelo freemium -- solo email y contraseña hasheada (RNF07), sin datos de pago, dirección ni información adicional -- en línea con los principios de minimización de datos de la Ley 25.326 de Protección de Datos Personales. Queda como trabajo pendiente para una eventual operación comercial real redactar una política de privacidad formal.

En cuanto a viabilidad técnica, el proyecto es enteramente de software: no requiere fabricación ni adquisición de hardware específico, dado que se ejecuta sobre infraestructura \emph{cloud} gestionada -- Render para el backend y Supabase para la base de datos PostgreSQL --, ambas con planes ya operativos y verificados en producción (Sección~\ref{sec:evidencia-cuantitativa}). Por lo tanto, este criterio también se da por cumplido sin requerir un análisis adicional de producción o compra de equipamiento.
